\documentclass[11pt, a4paper]{article}
\usepackage[margin=0.8in]{geometry}
\usepackage{enumitem}
\usepackage{titlesec}
\usepackage{xcolor}
\usepackage[UTF8]{ctex}
\usepackage{parskip}
\usepackage{hyperref}
\usepackage{array}
\usepackage{booktabs}
\usepackage{fontawesome5}
\usepackage{multicol}
\usepackage{ragged2e}
\usepackage{tabularx}

\definecolor{primary}{HTML}{2c3e50}  % 更深的蓝色，更专业
\definecolor{accent}{HTML}{3498db}   % 亮蓝色作为强调色
\definecolor{light}{HTML}{95a5a6}    % 灰色用于次要信息

% 设置超链接样式
\hypersetup{
    colorlinks=true,
    linkcolor=primary,
    filecolor=primary,
    urlcolor=accent,
    citecolor=primary,
}

% 自定义章节格式
\titleformat{\section}{
    \vspace{0.5em}\normalfont\Large\bfseries\color{primary}
}{}{0em}{}[\titlerule]

% 设置列表样式
\setlist[itemize]{
    leftmargin=*,
    topsep=2pt,
    parsep=0pt,
    partopsep=0pt,
    itemsep=2pt
}

% 移除段落缩进
\setlength{\parindent}{0pt}
\setlength{\parskip}{5pt}
\renewcommand{\arraystretch}{1.2}

% 自定义命令
\newcommand{\headeritem}[2]{\makebox[3.5em][l]{#1} #2}
\newcommand{\skillitem}[2]{\textbf{#1:} #2}
\newcommand{\dates}[1]{\hfill\textcolor{light}{\small #1}}  % 日期右对齐

\begin{document}
\pagestyle{empty}
\thispagestyle{empty}

% ========== 头部信息 ==========
\begin{center}
    {\Huge\bfseries\color{primary} 方知定} \\
    \vspace{0.3em}
    {\large\color{light} AI工程师 · 全栈开发者 · 5年经验} \\
    \vspace{0.8em}
    
    \begin{tabular}{@{}c@{}c@{}c@{}c@{}}
        \faPhone\ \href{tel:+8618170311560}{+86 181-7031-1560} & 
        \hspace{1em} | \hspace{1em} &
        \faEnvelope\ \href{mailto:1310347817@qq.com}{1310347817@qq.com} &
        \hspace{1em} | \hspace{1em} \\
        \faGithub\ \href{https://github.com/masterfzb}{github.com/masterfzb} &
        \hspace{1em} | \hspace{1em} &
        \faMapMarker\ 广州/深圳 &
        \hspace{1em} | \hspace{1em} \\
        \faGlobe\ \href{https://masterfzb.github.io/}{https://masterfzb.github.io/}
    \end{tabular}
\end{center}

\vspace{1em}

% ========== 专业摘要 ==========
\section*{\faUser\ 专业摘要}
致力于将前沿AI与工程技术转化为稳健的工业级解决方案。拥有5年全栈开发经验，专注金融科技、高并发系统与AI工程化。擅长通过逆向工程与系统重构解决复杂技术难题，具备将大语言模型等新兴技术落地到实际业务场景的成功经验。追求通过技术创新创造可量化的业务价值。

\vspace{-0.5em}

% ========== 工作经历 ==========
\section*{\faBriefcase\ 工作经历}

\subsection*{广发证券股份有限公司 \dates{2023.05 -- 2026.03}}
\textbf{金融大数据及后端研发工程师} \hfill \textit{广州}

\begin{itemize}
    \item \textbf{人工智能指数预测系统（核心项目）}
    \begin{itemize}
        \item 设计基于大语言模型的提示词工程体系，使用本地通义灵码自动化解析1734个指数文件，提取指标需求与算法
        \item 开发指标-数据库智能映射系统，通过LLM将数据需求转换为Oracle/MySQL查询，提升数据对接效率60%
        \item 构建Python量化计算引擎，根据指数算法自动化生成预测结果，并与真实市场数据进行验证比对
    \end{itemize}
    
    \item \textbf{证券业务高并发后端系统}
    \begin{itemize}
        \item 主导设计并开发基于Python+Django+Redis的证券后端系统，日均处理\textbf{百万级}数据请求，系统可用性99.9%
        \item 研发SQL动态生成工具，结合ORM与原生SQL优化复杂查询，关键报表响应时间减少40%
        \item 通过源码级分析修复内存泄漏等性能瓶颈，制定Git协作规范解决多团队并行开发问题
    \end{itemize}
\end{itemize}

\vspace{0.5em}

\subsection*{畅春互联科技有限公司 \dates{2021.10 -- 2023.02}}
\textbf{Python开发工程师} \hfill \textit{北京}

\begin{itemize}
    \item \textbf{政务信息智能采集与分类系统}
    \begin{itemize}
        \item 设计多线程分布式爬虫架构，突破反爬机制，实现目标网站全量数据采集
        \item 应用朴素贝叶斯算法实现网页信息自动分类，语义解析准确率达\textbf{90\%}，较人工提升70%
    \end{itemize}
    
    \item \textbf{移动安全自动化检测平台}
    \begin{itemize}
        \item 主导重构安卓动态检测工具，适配Android 13新权限模型并支持HarmonyOS跨平台兼容
        \item 参与研发业界首款鸿蒙系统Fuzzing测试工具，通过指令变异引擎发现多个系统级漏洞
    \end{itemize}
\end{itemize}

\vspace{0.5em}

\subsection*{荣耀 \dates{2021.04 -- 2021.08}}
\textbf{Python测试开发工程师} \hfill \textit{北京}

\begin{itemize}
    \item 开发自动化测试脚本，用于手机及智能手表功耗与性能测试，根据需求动态调整测试参数
    \item 自动生成包含性能指标、能耗分析在内的详细测试报告，提升测试效率30%
\end{itemize}

\vspace{0.5em}

\subsection*{上海宝信软件股份有限公司 \dates{2019.01 -- 2020.02}}
\textbf{组态软件工程师} \hfill \textit{上海}

\begin{itemize}
    \item 参与苏州地铁3号线、乌鲁木齐地铁TCC等多个大型轨道交通综合监控系统项目
    \item 负责上位机监控画面组态开发、现场部署与调试，利用自动化脚本提升工作效率30%
\end{itemize}

% ========== 技术能力 ==========
\section*{\faCogs\ 技术能力}

\begin{multicols}{2}
\begin{itemize}
    \item \textbf{编程语言}: Python, SQL, C++, Bash
    \item \textbf{后端开发}: Django, RESTful API, Redis, 高并发设计
    \item \textbf{AI/机器学习算法}: 大模型应用, 提示词工程, 机器学习算法, 生成模型
    \item \textbf{数据工程}: Scrapy, 数据挖掘, 数据库优化, 大数据处理
    \item \textbf{DevOps}: Git, Linux, Docker, CI/CD, 自动化部署
    \item \textbf{系统工程}: 架构设计, 性能优化, 系统重构, 逆向工程
    \item \textbf{智能体工程}: openclaw部署及工程运维,skill制作,ollama模型部署,AI全栈开发
\end{itemize}
\end{multicols}

% ========== 项目亮点 ==========
\section*{\faStar\ 项目亮点}

\begin{itemize}
    \item \textbf{大模型金融工程化}: 成功将LLM技术应用于金融指数分析全流程，从文档解析、需求提炼到计算验证，构建自动化智能分析管道
    \item \textbf{复杂系统重构}: 多次接手遗留系统，通过逆向工程与动态分析定位核心缺陷，设计可持续改进方案解决历史性技术债务
    \item \textbf{全流程开发能力}: 独立负责项目从需求分析、架构设计、核心开发到生产部署与运维的全周期
    \item \textbf{前沿技术探索}: 持续实践AI前沿技术，成功部署OpenClaw、Ollama、GPT-SoVITS等开源框架，开发分角色有声书制作工具
\end{itemize}

% ========== 教育背景 ==========
\section*{\faGraduationCap\ 教育背景}

\subsection*{上海海事大学 \dates{2014.09 -- 2018.06}}
\textbf{工学学士 | 电气工程及其自动化}

\begin{itemize}
    \item \textbf{主修课程}: 深度学习, 计算机系统结构, C语言程序设计, 电路原理, 数字电子技术
    \item \textbf{学业表现}: GPA 3.66/5 (专业前50\%), 大学物理竞赛上海赛区二等奖
    \item \textbf{核心能力}: 全国计算机二级, 英语六级, 熟练掌握Office, LaTeX等工具
\end{itemize}

% ========== 实习经历 ==========
\section*{\faLaptopCode\ 实习经历}

\subsection*{INTAMSYS（远铸智能） \dates{2018.04 -- 2018.06}}
\textbf{电气工程实习生}

\begin{itemize}
    \item 参与工业级3D打印设备的电气装配、线路检查与设备调试
    \item 协助更新与维护设备电气原理图及技术文档
\end{itemize}

% ========== 页脚 ==========
\vfill
\begin{center}
    \small\textcolor{light}{最后更新: 2026年3月 · 《林无静木，川无停流》}
\end{center}

\end{document}